\label{sec:methodology}

\subsection{The \textit{Karcher} Doctrine in Detail}

Justice Brennan's majority opinion in \textit{Karcher v.\ Daggett} is the locus
classicus for incumbency as a redistricting criterion, and it deserves careful attention
because its language is frequently misread in redistricting proceedings.

The holding of \textit{Karcher} was that New Jersey's congressional map, which contained
population deviations of 0.6984\%, was unconstitutional because the state failed to
demonstrate that those deviations were necessary to achieve a legitimate state objective.
Justice Brennan's majority accepted that the deviations were small in absolute terms
but held that the Equal Population principle demanded justification for any deviation
from mathematical equality in congressional apportionment~\citep{karcher1983}.

The language on legitimate objectives appeared in the context of explaining what kinds
of justifications would suffice to uphold a map with minor deviations:

\begin{quote}
``Any number of consistently applied legislative policies might justify some
variance, including, for example, making districts compact, respecting municipal
boundaries, preserving the cores of prior districts, and avoiding contests between
incumbents.''
\end{quote}

Three features of this passage are critical for the incumbency objection. First, the
passage is framed hypothetically: it describes what ``might'' justify deviations, not
what must be done. Second, it is explicitly a non-exhaustive list (``for example''),
indicating that the Court was not enumerating mandatory criteria. Third, the passage
appears in a section explaining why the New Jersey map \textit{failed} the test---the
state did not demonstrate that the listed or other objectives justified its deviations.
The Court was not endorsing the criteria; it was explaining what a successful justification
would need to show.

The misreading of \textit{Karcher} in incumbency litigation typically proceeds as follows:
the \textit{Karcher} dicta identifies incumbency protection as a legitimate objective;
a legitimate objective is a required consideration; therefore, any redistricting plan
that does not consider incumbency fails the \textit{Karcher} test. Each step in this
syllogism is incorrect. ``Legitimate objective'' means an objective that is not
constitutionally prohibited, not one that is constitutionally required. Redistricting
bodies are free to pursue or not pursue any legitimate objective, so long as they
satisfy the constitutional requirements of equal population, non-dilution of minority
voting strength, and contiguity.

\subsection{The \textit{Cox v.\ Larios} Limitation}

\textit{Cox v.\ Larios} provides the clearest statement of the limits of incumbency
as a justification for population deviations. The three-judge district court in the
Northern District of Georgia found that Georgia's state legislative maps contained
population deviations of up to 9.98\% in the state Senate and 4.98\% in the state House,
justified in part by the legislature's desire to protect incumbent legislators. The
court held that using incumbency protection as a primary justification for deviations
of this magnitude was unconstitutional: the deviations were too large to be justified
by any legitimate interest, including incumbency protection~\citep{cox2004}.

The Supreme Court's summary affirmance of \textit{Cox v.\ Larios} confirmed that
the district court's analysis was correct, at least as applied to the facts. The holding
can be read as establishing a rough proportionality constraint: small deviations may
be justified by legitimate interests including incumbency; large deviations cannot.
For congressional redistricting, where any deviation from mathematical equality must
be justified~\citep{karcher1983}, the \textit{Cox} holding implies that incumbency
protection cannot justify any deviation that is not the \textit{necessary} result of
pursuing the incumbent-protection objective---a high standard that enacted maps
frequently fail to meet.

\subsection{Algorithmic Redistricting and Population Deviation}

The \textsc{bisect} algorithm produces congressional maps that satisfy the
equal-population constraint without any incumbency-related deviation. By operating
on census tracts rather than precincts, and by explicitly balancing population at
each bisection step, the algorithm produces maps where the maximum population deviation
across districts is well below 1\% in all three states studied. This performance is
not dependent on incumbency protection---the algorithm achieves near-perfect population
equality purely through its population-balance optimisation.

This means that the \textit{Karcher}/\textit{Cox} analysis is largely irrelevant to
algorithmic redistricting: the algorithm does not introduce population deviations that
require justification by incumbency or any other criterion. The \textit{Karcher}
line of cases addresses the justification for deviations that the algorithm does not
produce. The incumbency objection, framed as a \textit{Karcher} challenge, therefore
fails at the threshold: there are no deviations to justify.

\subsection{Three Uses of Algorithmic Incumbency Analysis in Litigation}

The I.1--I.3 empirical results can be deployed in redistricting litigation in three
distinct ways.

\paragraph{Establishing non-randomness of enacted pairing rates.}
I.1's binomial test demonstrates that enacted map pairing rates are statistically
inconsistent with random redistricting. This is evidence that mapmakers actively used
incumbent address data to minimise pairings---a redistricting choice, not a geographic
necessity. In litigation challenging the enacted map, this evidence supports the
inference that the map was drawn with at least some degree of strategic incumbency
protection.

\paragraph{Providing a neutral reference distribution.}
The \textsc{bisect} map provides a neutral reference point for all three incumbency
metrics: pairing rate, safe-seat fraction, and open-seat count. Courts can use this
reference distribution to assess whether the enacted map's incumbency outcomes are
within the range of what geographic redistricting would produce, or whether they
represent deliberate departures from geographic chance. A large positive incumbency
protection premium (as documented in I.0--I.3) supports the inference of strategic
mapmaking.

\paragraph{Quantifying the intentional incumbency protection premium.}
The IPP metric (I.1) provides a single number that quantifies the departure from
geographic neutrality on incumbent pairing. Courts evaluating whether an enacted map's
use of incumbency was proportionate to its legitimate objectives can use the IPP to
assess the degree of the departure: a large IPP (enacted pairing rate far below the
geographic baseline) suggests that incumbency was not merely considered but was a
dominant criterion that displaced other redistricting objectives.
